\section{Methods}\label{sec:methods}

% TODO: describe the methodology used to construct whatever this paper needs
% before the audits in \Cref{sec:primary-results,sec:secondary-results} apply.
% Stages typical in this series are listed below; rename, drop, or add to fit
% the actual route used by this paper.

\subsection{Method overview}\label{subsec:method-overview}

% TODO: one-paragraph overview naming the route as a whole, citing the
% inherited frameworks it composes, and summarizing what the route produces.

\subsection{Prerequisites}\label{subsec:method-prerequisites}

% TODO: list the upstream constructions this method takes as given.

\subsection{Source-of-truth policy}\label{subsec:source-policy}

% TODO: state the source-policy profile (which source classes are admissible
% for which role claims under which instrument).

\subsection{Claim-boundary discipline}\label{subsec:claim-boundary-discipline}

% TODO: state how this paper separates positive claims, retained compatibility
% claims, and explicit nonclaims.  Point forward to the controls that enforce
% the boundary.

\subsection{Candidate inputs and admission gate}\label{subsec:candidate-admission}

% TODO: name the candidate objects entering the construction and the gate that
% admits them.

\subsection{Construction step}\label{subsec:construction-step}

% TODO: the principal construction step (e.g., exactification, completion,
% promotion-then-exactification, etc.).

\subsection{Native structure of the constructed objects}\label{subsec:native-structure}

% TODO: describe the native structure of the constructed objects (state space,
% laws, splits, invariance records).

\subsection{Irreducibility certificate}\label{subsec:irreducibility-certificate}

% TODO: record the construction-side irreducibility / non-factorization
% certificate that the rest of the paper depends on.

\subsection{Constructed-object table}\label{subsec:constructed-object-table}

% TODO: drop in a summary table of the constructed objects.
% % TEMPLATE: Table 2.  Typical use in this series is the constructed-object
% summary in \Cref{subsec:constructed-object-table}.  Replace caption, label,
% column headers, and rows.
\begin{table}[tbp]
  \centering
  \caption{TODO: caption.}
  \label{tab:template-02}
  \small
  \begin{tabularx}{\textwidth}{@{}>{\raggedright\arraybackslash}p{0.20\textwidth}>{\raggedright\arraybackslash}X>{\raggedright\arraybackslash}p{0.17\textwidth}>{\raggedright\arraybackslash}p{0.16\textwidth}>{\raggedright\arraybackslash}p{0.18\textwidth}@{}}
    \toprule
    Column 1 & Column 2 & Column 3 & Column 4 & Column 5 \\
    \midrule
    TODO row
      & TODO
      & TODO
      & TODO
      & TODO \\
    \bottomrule
  \end{tabularx}
\end{table}

